% Part: sets-functions-relations
% Chapter: infinite
% Section: dedekinds-proof
%
\documentclass[../../../include/open-logic-section]{subfiles}

\begin{document}

\olfileid{sfr}{infinite}{dedekindsproof}

\olsection[Dedekinds ``bewijs'']{Dedekinds ``bewijs'' dat er een
oneindige verzameling bestaat}

In dit hoofdstuk hebben we de natuurlijke getallen met verzamelingen
beschreven. Daarbij gebruikten we Dedekindalgebra's. In
\olref[arith][ref]{sec} bespraken we onder meer wat de aritmetisering
van de analyse filosofisch betekent: de analyse wordt in rekenkundige
termen opgebouwd. Nu moeten we bekijken wat het resultaat van dit
hoofdstuk betekent.

In heel \olref[sfr][arith][]{chap} gingen we uit van de natuurlijke
getallen en lieten we precies zien hoe we daarmee de gehele, rationale
en reële getallen konden opbouwen. In dit hoofdstuk hebben we de
natuurlijke getallen zelf niet op die manier opgebouwd. We hebben
alleen laten zien dat we \emph{van elke Dedekind-oneindige verzameling}
een verzameling kunnen definiëren die precies de eigenschappen heeft die
we van~$\Nat$ verlangen.

We kunnen dus niet beweren dat we een metafysische vraag hebben
beantwoord, bijvoorbeeld: \emph{welke objecten zijn de natuurlijke
getallen}? Dat is maar goed ook. In
\olref[sfr][arith][ref]{sec} zeiden we immers dat we de definitie
van~$\Real$ als de verzameling van Dedekindsneden niet als een
\emph{ontdekking} moeten zien. Het is een handige afspraak. Waar het om
gaat, is dat Dedekindsneden de belangrijkste wiskundige eigenschappen
van de reële getallen hebben. Hier geldt hetzelfde: \emph{elke}
Dedekindalgebra heeft de belangrijkste wiskundige eigenschappen van de
natuurlijke getallen. (Dedekind maakte dit extra duidelijk door te
bewijzen dat alle Dedekindalgebra's \emph{isomorf} zijn
(\citeyear[stellingen 132--3]{Dedekind1888}). Geen wonder dat veel
hedendaagse ``structuralisten'' Dedekind daarom als voorloper noemen.)

We hebben ook laten zien hoe we de theorie van de natuurlijke getallen
kunnen onderbrengen in een eenvoudige, na\"ieve verzamelingenleer. Die
leer is nog vrij informeel. Maar ze lijkt de natuurlijke getallen niet
vooraf als gegeven aan te nemen. Misschien komen we zo dichter bij
Dedekinds eigen ambitieuze plan. Hij legde het zo uit:
\begin{quote}
	Als iets bewezen kan worden, mogen we het in de wetenschap niet zonder
	bewijs geloven. Dat lijkt een redelijke eis. Toch vind ik dat zelfs de
	nieuwste manieren om de basis van de eenvoudigste wetenschap te leggen
	nog niet aan die eis voldoen. Daarmee bedoel ik het deel van de logica
	dat over de getallentheorie gaat. Als ik rekenkunde (algebra, analyse)
	slechts een deel van de logica noem, bedoel ik dat het begrip getal
	volgens mij helemaal losstaat van onze begrippen of intuïties van
	ruimte en tijd---ik zie het juist als iets dat rechtstreeks voortkomt
	uit de zuivere wetten van het denken.
	\citep[voorwoord]{Dedekind1888}
\end{quote}
Dedekinds gedurfde idee is dus dit. We hebben net laten zien hoe we de
natuurlijke getallen alleen met (na\"ieve) verzamelingenleer kunnen
opbouwen. In \olref[sfr][arith][]{chap} zagen we hoe we de reële
getallen kunnen opbouwen met de natuurlijke getallen als uitgangspunt
en met wat verzamelingenleer. Misschien blijken ``rekenkunde (algebra,
analyse)'' dus toch ``slechts een deel van de logica'' te zijn (in de
ruimere betekenis die Dedekind aan het woord ``logica'' geeft).

Dat is het idee. Maar wacht even. Om een Dedekindalgebra te
bouwen---onze vervanger voor de natuurlijke getallen---moet er wel een
Dedekind-oneindige verzameling bestaan. (Kijk maar terug naar
\olref[sfr][infinite][dedekind]{thm:DedekindInfiniteAlgebra}.) Als we
het bestaan van zo'n verzameling niet met ``logica'' of ``de zuivere
wetten van het denken'' kunnen bewijzen, loopt het plan vast.

\emph{Kunnen} we met ``de zuivere wetten van het denken'' dan bewijzen
dat er een Dedekind-oneindige verzameling bestaat? Dedekind probeerde
het zo:
\begin{quote}
  Mijn eigen wereld van gedachten, dat wil zeggen het geheel~$S$ van alle
  dingen waarover ik kan denken, is oneindig. Als~$s$ voor een element
  van~$S$ staat, dan is de gedachte~$s'$ dat ik over~$s$ kan denken zelf
  ook een element van~$S$. Als we die gedachte zien als een
  beeld~$\phi(s)$ van het element~$s$, dan is \dots~$S$
  [Dedekind-]oneindig, en dat moest worden bewezen.
	\citep[\S66]{Dedekind1888}
\end{quote}
Dit is nogal verbazingwekkend in een boek dat verder grotendeels uit
zeer precieze wiskundige bewijzen bestaat. Hier vallen twee dingen op.

Ten eerste: dit ``bewijs'' zouden we tegenwoordig nauwelijks wiskundig
noemen. Het gaat over psychologische objecten, namelijk gedachten, en
dan ook nog over gedachten die alleen \emph{mogelijk} zijn.

Ten tweede: met Dedekindalgebra's zoals wij die hebben uitgelegd, zit er
ook een duidelijk technisch gat in dit ``bewijs''. Als Dedekinds
redenering klopt, bewijst ze alleen dat er oneindig veel dingen zijn
(om precies te zijn: oneindig veel gedachten). Maar Dedekind moet ons
ook een reden geven om~$S$ te zien als één \emph{verzameling} met
oneindig veel elementen. In plaats daarvan kan~$S$ gewoon staan voor
\emph{een aantal dingen} (meervoud). Die hoeven samen nog geen
verzameling te vormen.

Dedekind zag dit gat niet. Dat kan betekenen dat zijn woord
``totaliteit'' niet precies hetzelfde betekende als wat \emph{wij} met
``verzameling'' bedoelen.\footnote{Er zijn inderdaad andere redenen om
dat te denken; zie \citet[p.~23]{Potter2004}.} Dat zou niet erg
verrassend zijn. In de afgelopen twee hoofdstukken hebben we alles op
een na\"ieve manier gedaan: we maakten een ``constructie'' van de
natuurlijke getallen en bouwden daaruit de gehele, reële en rationale
getallen op. Telkens als we een bepaalde verzameling nodig hadden,
namen we gewoon aan dat die bestond. We legden geen algemene regels
vast die precies zeggen welke verzamelingen bestaan en wanneer. Maar we
weten dat \emph{enkele} algemene regels nodig zijn. Anders krijgen we
de paradox van Russell.

Het is tijd om niet langer zo na\"ief te werk te gaan.
\end{document}